% Все презентации курса компилируются только pdflatex.
\usepackage{iftex}
\RequirePDFTeX
\usepackage{cmap}
\usepackage[T2A]{fontenc}
\usepackage[utf8]{inputenc}
\usepackage[english,russian]{babel}
\usepackage{PTSerif}
\usepackage{amsmath,amsfonts,amssymb}
\usepackage{graphicx,microtype,anyfontsize,tikz}
\usepackage{ragged2e}
\usetikzlibrary{arrows.meta,calc}

\usetheme{default}
\usefonttheme{serif}
\DeclareMathOperator{\sen}{sen}
\bibliographystyle{apalike}

% Общие сведения; номер лекции, её название и дата задаются в N.tex.
\newcommand{\coursename}{Введение в маневрирование КА}
\author{Чиняев Владимир Николаевич}
\institute{Московский физико-технический институт}

% Книжное оформление с голубым акцентом.
\definecolor{coursepaper}{HTML}{FBF9F3}
\definecolor{courseink}{HTML}{302B28}
\definecolor{courseaccent}{HTML}{187AA5}
\definecolor{coursemuted}{HTML}{6D6560}
\definecolor{courserule}{HTML}{A9CADA}
\definecolor{coursetint}{HTML}{EDF4F6}
\tikzset{
  course diagram/.style={font=\fontsize{9.5}{11.5}\selectfont,
    line cap=round,line join=round,>=Latex},
  orbit line/.style={draw=courseaccent,line width=1pt},
  aux line/.style={draw=courserule,line width=.5pt},
  vector line/.style={draw=courseink,line width=.65pt,->},
  angle line/.style={draw=courseaccent,line width=.8pt,->},
  orbit point/.style={circle,fill=courseaccent,inner sep=1.7pt}
}
\setbeamercolor{normal text}{fg=courseink,bg=coursepaper}
\setbeamercolor{background canvas}{bg=coursepaper}
\setbeamercolor{structure}{fg=courseaccent}
\setbeamercolor{alerted text}{fg=courseaccent}
\setbeamercolor{frametitle}{fg=courseink,bg=coursepaper}
\setbeamercolor{footline}{fg=coursemuted,bg=coursepaper}
\setbeamercolor{caption name}{fg=courseaccent}

% Основной текст — 11/14 pt; подписи к рисункам — 9.5/11.5 pt.
\let\coursenormalsize\normalsize
\renewcommand{\normalsize}{\coursenormalsize\fontsize{11}{14}\selectfont}
\setbeamerfont{normal text}{size*={11pt}{14pt}}
\setbeamerfont{frametitle}{size*={20pt}{23pt},series=\mdseries}
\setbeamerfont{title}{size*={27pt}{31pt},series=\mdseries}
\setbeamerfont{subtitle}{size*={12.8pt}{15pt},series=\mdseries}
\setbeamerfont{author}{size*={10.8pt}{13pt}}
\setbeamerfont{institute}{size*={9.2pt}{11pt}}
\setbeamerfont{date}{size*={9pt}{11pt}}
\setbeamerfont{footline}{size*={6.8pt}{8pt}}
\setbeamerfont{caption}{size*={9.5pt}{11.5pt}}
\setbeamerfont{itemize/enumerate body}{size*={11pt}{14pt}}
\setbeamerfont{itemize/enumerate subbody}{size*={11pt}{14pt}}
\setbeamerfont{itemize/enumerate subsubbody}{size*={11pt}{14pt}}
\AtBeginDocument{\normalsize}

\setbeamersize{text margin left=6mm,text margin right=6mm}
\setbeamertemplate{navigation symbols}{}
\setbeamertemplate{headline}{}
\setbeamertemplate{caption}[numbered]
\setbeamertemplate{itemize items}[circle]

% Стандартные блоки Beamer позволяют группировать текст и формулы.
\setbeamertemplate{blocks}[default]
\setbeamerfont{block title}{size*={11pt}{14pt},series=\bfseries}
\setbeamerfont{block body}{size*={11pt}{14pt}}
\setbeamercolor{block title}{fg=courseaccent,bg=coursepaper}
\setbeamercolor{block body}{fg=courseink,bg=coursetint}
\setbeamercolor{block title alerted}{parent=block title}
\setbeamercolor{block body alerted}{parent=block body}
\setbeamercolor{block title example}{parent=block title}
\setbeamercolor{block body example}{parent=block body}

% Компактный заголовок; его высота учитывается при размещении содержания.
\setbeamertemplate{frametitle}{%
  \vskip4.5mm%
  \begin{beamercolorbox}[wd=\textwidth,sep=0pt]{frametitle}
    \usebeamerfont{frametitle}\strut\insertframetitle\strut\par
    \vskip2mm%
    {\color{courseaccent}\hrule height .55pt}
  \end{beamercolorbox}%
}

% Нижний колонтитул: курс слева, лекция и номер слайда справа.
\setbeamertemplate{footline}{%
  \begin{beamercolorbox}[wd=\paperwidth,ht=5.5mm,dp=0pt]{footline}
    \begin{tikzpicture}[remember picture,overlay,x=1mm,y=1mm]
      \begin{scope}[shift={(current page.south west)}]
        \draw[courserule,line width=.3pt] (6,5.5)--(154,5.5);
        \node[text=coursemuted,anchor=west,inner sep=0pt,
          font=\usebeamerfont{footline}] at (6,2.8) {\coursename};
        \node[text=coursemuted,anchor=east,inner sep=0pt,
          font=\usebeamerfont{footline}] at (154,2.8)
          {Лекция \arabic{lecture}\quad\textbullet\quad\insertframenumber};
      \end{scope}
    \end{tikzpicture}%
  \end{beamercolorbox}%
}

% Титульник рассчитан на 16:9 (160 × 90 мм), без нижнего колонтитула.
\setbeamertemplate{title page}{%
  \begin{tikzpicture}[remember picture,overlay,x=1mm,y=1mm]
    \begin{scope}[shift={(current page.south west)}]
      \node[text=coursemuted,anchor=north,font=\usebeamerfont{institute}]
        at (80,82) {\insertinstitute};
      \draw[courseaccent,line width=.55pt] (71,71)--(89,71);
      \node[text=courseink,align=center,text width=125mm,
        font=\usebeamerfont{title}] at (80,53)
        {\hyphenpenalty=10000\coursename\par};
      \node[text=courseaccent,font=\usebeamerfont{subtitle}]
        at (80,32) {Лекция \arabic{lecture}. \inserttitle};
      \node[text=courseink,font=\usebeamerfont{author}]
        at (80,19) {\insertauthor};
      \node[text=coursemuted,font=\usebeamerfont{date}]
        at (80,9) {\insertdate};
    \end{scope}
  \end{tikzpicture}%
}
